%*******************************************************************************
%*********************************** First Chapter *****************************
%*******************************************************************************

\chapter{Introduzione} 
\label{cap:1}

Il presente capitolo si pone l'obiettivo di delineare il perimetro concettuale e scientifico entro cui si sviluppa la tesi, focalizzandosi sul confronto critico tra due delle metodologie più promettenti e radicalmente diverse del panorama contemporaneo del Machine Learning: la \textbf{Knowledge Distillation} e il \textbf{Reflection Tuning}. Sebbene entrambe le tecniche convergano verso l'obiettivo comune di ottenere modelli linguistici più performanti e sostenibili, esse rappresentano filosofie di ottimizzazione distinte. Per comprendere la portata di questo confronto, è tuttavia indispensabile ricostruire il percorso evolutivo che ha portato la disciplina a scontrarsi con i limiti della scala dimensionale, rendendo necessario un cambio di paradigma verso tecniche di raffinamento procedurale.

Sino a pochi anni fa, l'evoluzione dell'apprendimento automatico è stata guidata da una fiducia quasi incondizionata nella crescita dimensionale delle architetture. Prima che concetti come la distillazione o l'auto-riflessione diventassero centrali, la ricerca era focalizzata sul superamento dei limiti del Machine Learning classico. In quella fase embrionale, il passaggio cruciale è stato l'abbandono del \textit{feature engineering} manuale — dove l'uomo istruiva la macchina su quali caratteristiche osservare — in favore di una rappresentazione automatica e gerarchica dei dati tramite reti neurali profonde.

Il \textbf{Deep Learning}, affermatosi definitivamente con il successo di AlexNet nella competizione ImageNet del 2012~\cite{krizhevsky2012imagenet}, ha dimostrato che le reti neurali con molteplici strati nascosti erano capaci di apprendere autonomamente rappresentazioni sempre più astratte e significative dei dati grezzi. Le prime architetture si concentravano prevalentemente su dati strutturati e su modalità specifiche, come le \textit{Convolutional Neural Networks} (CNN) per la visione artificiale e le \textit{Recurrent Neural Networks} (RNN) per le sequenze temporali. Queste ultime, nella variante \textit{Long Short-Term Memory} (LSTM)~\cite{hochreiter1997long}, avevano l'ambizione di modellare le dipendenze a lungo termine nel linguaggio naturale, ma soffrivano di limitazioni intrinseche: l'addestramento era intrinsecamente sequenziale, la memorizzazione di contesti molto lunghi rimaneva problematica e la parallelizzazione su hardware moderno risultava inefficiente.

Tuttavia, nonostante il salto tecnologico, l'addestramento rimaneva ancorato a una visione puramente statistica e "passiva". Il modello apprendeva attraverso la minimizzazione di una funzione di perdita basata su etichette rigide (\textit{hard labels}), cercando di mappare gli input sugli output in modo diretto. Questo approccio, seppur efficace per task di classificazione specifici, operava come una sorta di memoria associativa complessa: il modello non catturava la semantica profonda o le sfumature di incertezza intrinseche nei dati, limitando la sua capacità di generalizzare a contesti mai visti durante il training. L'intelligenza era, in questa fase, un riflesso statistico della frequenza dei dati di input.

La svolta decisiva è arrivata nel 2017 con la pubblicazione di \textit{"Attention Is All You Need"}~\cite{vaswani2017attention} da parte di ricercatori di Google Brain. Questo lavoro ha introdotto l'architettura \textbf{Transformer}, che ha soppiantato le reti ricorrenti come paradigma dominante per il trattamento delle sequenze. Il meccanismo fondante di questa architettura è l'\textit{self-attention}: anziché processare i token di una sequenza in modo strettamente seriale, il Transformer calcola, per ogni elemento della sequenza, un peso di attenzione verso tutti gli altri elementi contemporaneamente. Questo consente di catturare dipendenze a lungo raggio in modo diretto ed efficiente, indipendentemente dalla distanza tra i token nella sequenza.

Il secondo vantaggio fondamentale dei Transformer risiede nella loro intrinseca \textbf{parallelizzabilità}. A differenza delle RNN, dove il calcolo al passo $t$ dipende dal risultato al passo $t-1$, tutte le operazioni di attenzione possono essere eseguite in parallelo su unità di calcolo come le GPU. Questo ha reso possibile addestrare architetture su scale prima impensabili, aprendo la strada all'era dei Large Language Models (LLM).

I primi modelli a sfruttare appieno questa architettura per il linguaggio naturale sono stati GPT~\cite{radford2018improving} di OpenAI, basato su un decoder causale pre-addestrato con l'obiettivo di predizione del token successivo, e BERT~\cite{devlin2018bert} di Google, che ha introdotto il pre-addestramento bidirezionale tramite \textit{Masked Language Modeling}. Entrambi hanno stabilito il paradigma del \textbf{pre-training e fine-tuning}: addestrare un modello generico su enormi quantità di testo non etichettato e poi specializzarlo su task specifici con dati limitati. Questo approccio ha dimostrato un'eccezionale capacità di trasferimento della conoscenza, consentendo ai modelli di affrontare con pochi esempi (\textit{few-shot}) o addirittura senza esempi (\textit{zero-shot}) task per i quali non erano stati esplicitamente addestrati.

La traiettoria evolutiva successiva è stata segnata dal modello GPT-3~\cite{brown2020language}, con i suoi 175 miliardi di parametri, che ha rappresentato un punto di discontinuità qualitativa. GPT-3 ha mostrato che, al di sopra di certe soglie dimensionali, i modelli linguistici sviluppano \textit{capacità emergenti}: abilità che non erano presenti nei modelli più piccoli e che non erano state esplicitamente insegnate, come il ragionamento analogico, la scrittura di codice o la risoluzione di problemi matematici elementari. Si è così delineata la classe dei \textbf{Large Language Models}: modelli con decine o centinaia di miliardi di parametri, addestrati su corpus testuali di scala planetaria, capaci di interagire in linguaggio naturale con una fluidità e una versatilità senza precedenti.

\section{Il dilemma delle Scaling Laws}
\label{sec:scaling}

Con l'avvento dell'architettura Transformer e l'esplosione dei Large Language Models, il settore è entrato in una fase caratterizzata dalle cosiddette \textit{Scaling Laws}~\cite{kaplan2020scaling}. L'idea prevalente, supportata da evidenze empiriche inizialmente indiscutibili, era che la semplice aggiunta di parametri, unita a volumi di dati sempre più massicci e potenza di calcolo crescente, avrebbe naturalmente risolto ogni lacuna cognitiva delle macchine. Kaplan et al. hanno formalizzato questa intuizione dimostrando che le performance dei modelli linguistici seguono leggi di potenza prevedibili rispetto a tre variabili: numero di parametri, dimensione del dataset e budget computazionale. Questa corsa all'iper-parametrizzazione ha effettivamente prodotto risultati sbalorditivi, portando alla comparsa di capacità emergenti che sembravano mimare il ragionamento umano.

Tuttavia, questo successo ha contemporaneamente generato una crisi di sostenibilità su più fronti, che costituisce il punto di partenza della presente analisi. Il primo e più evidente ostacolo è di natura infrastrutturale: l'esecuzione di modelli con centinaia di miliardi di parametri richiede una potenza di calcolo e un consumo energetico non più trascurabili. L'addestramento di GPT-4, secondo stime pubblicamente disponibili, avrebbe richiesto decine di migliaia di GPU specializzate per mesi, con un costo stimato nell'ordine delle decine di milioni di dollari~\cite{patel2023gpt4}. Anche la sola inferenza — la semplice esecuzione del modello per rispondere a una query — richiede hardware dedicato ad alte prestazioni, rendendo l'IA d'avanguardia un privilegio di pochissimi attori industriali con risorse straordinarie ed escludendo di fatto ricercatori indipendenti, università e piccole imprese. Si configura così un rischio concreto di concentrazione tecnologica che pone interrogativi fondamentali sulla democratizzazione dell'accesso all'intelligenza artificiale.

A questo si affianca un secondo ostacolo pratico, spesso sottovalutato nel dibattito pubblico: la latenza di inferenza. Un'applicazione interattiva — un assistente virtuale, un sistema di raccomandazione in tempo reale, uno strumento di supporto medico — necessita di risposte nell'ordine dei millisecondi. I modelli di grandi dimensioni, anche quando ospitati su infrastrutture dedicate, introducono latenze incompatibili con molti scenari d'uso reali. Il dispiegamento su dispositivi \textit{edge} — smartphone, dispositivi IoT, sistemi embedded — è poi completamente precluso quando il modello richiede diversi gigabyte solo per caricare i propri pesi in memoria, rendendo di fatto impossibile qualsiasi forma di intelligenza artificiale locale e privata per l'utente finale.

Al di là dei limiti infrastrutturali, si è palesata con altrettanta chiarezza la fragilità logica intrinseca di questi sistemi. Pur essendo capaci di generare testi fluenti, i modelli risultano spesso inclini alle \textbf{allucinazioni}~\cite{ji2023survey}: la produzione di affermazioni sintatticamente coerenti ma fattualmente false o logicamente inconsistenti. Questo accade perché il processo di generazione rimane greedy o stocastico — il modello seleziona il token successivo sulla base della distribuzione di probabilità appresa, senza una reale validazione della coerenza logica complessiva dell'output che si sta costruendo. La mancanza di un meccanismo di \textit{self-checking} intrinseco limita profondamente l'affidabilità di questi sistemi in contesti ad alto rischio, come la diagnosi medica, la consulenza legale o i sistemi di controllo autonomo, dove un'allucinazione non è un mero inconveniente ma può tradursi in conseguenze concrete e dannose.

Infine, studi più recenti hanno messo in discussione la stessa praticabilità della crescita parametrica come strategia a lungo termine. Il lavoro di Hoffmann et al. su Chinchilla~\cite{hoffmann2022training} ha dimostrato che molti modelli della generazione precedente erano sovra-parametrizzati rispetto alla quantità di dati di addestramento disponibili: la strategia ottimale prevede un bilanciamento preciso tra numero di parametri e volume di token di training. Ciò implica che la comunità scientifica si trova a fare i conti non solo con un limite computazionale ed economico, ma con un limite fisico nella disponibilità di dati testuali di qualità su scala internet. La semplice aggiunta di parametri, senza un corrispondente aumento dei dati, produce rendimenti marginali decrescenti, segnando di fatto l'esaurimento della strategia del gigantismo come via maestra verso modelli più capaci.

In questo scenario di gigantismo tecnologico insostenibile, la comunità scientifica ha iniziato a interrogarsi su come rendere l'intelligenza artificiale non solo più potente, ma più \textbf{intelligente ed efficiente}. La domanda centrale si è spostata: anziché chiedersi come costruire modelli sempre più grandi, ci si è chiesti come estrarre il massimo valore da modelli di dimensioni gestibili. Da questa ricerca di efficienza scaturisce il confronto al centro di questo lavoro.

La \textbf{Knowledge Distillation} è emersa come la prima risposta strutturata al problema della scala~\cite{hinton2015distilling}. Essa propone un ribaltamento del concetto di addestramento: invece di istruire un modello piccolo partendo da zero su dati grezzi, si utilizza un modello più vasto e già esperto (\textit{Teacher}) per guidarne uno più snello (\textit{Student}). La forza della distillazione risiede nella cosiddetta "conoscenza oscura" (\textit{Dark Knowledge}). Il docente non trasmette solo la risposta corretta, ma anche le distribuzioni di probabilità sulle risposte errate, rivelando allo studente le relazioni semantiche e le gerarchie di similarità che il modello grande ha faticosamente appreso. È una strategia di \textbf{efficienza per imitazione}: lo studente diventa un'ombra computazionalmente leggera del maestro.

Tuttavia, la compressione strutturale della distillazione eredita spesso i limiti logici del docente. Un modello studente distillato tende a replicare non solo le capacità del teacher, ma anche i suoi pattern di errore e le sue allucinazioni caratteristiche. Per questo motivo, la tesi introduce e contrappone il \textbf{Reflection Tuning}~\cite{ze2023reflection}, una tecnica che rappresenta l'avanguardia del raffinamento qualitativo. Se la distillazione guarda al trasferimento \textit{esterno} del sapere, il Reflection Tuning punta allo sviluppo di una competenza \textit{interna} attraverso l'introspezione. Esso trae ispirazione dai processi metacognitivi umani — il cosiddetto Sistema 2 di Kahneman~\cite{kahneman2011thinking}, ovvero il pensiero lento, deliberato e critico — addestrando il modello a non fornire la prima risposta utile, ma a sottoporre il proprio output a un ciclo strutturato di analisi e revisione. Questo passaggio da una generazione \textit{one-shot} a un processo iterativo di auto-correzione permette di abbattere le allucinazioni, spostando l'enfasi dalla quantità di parametri alla qualità del processo di inferenza. Si tratta di una strategia di \textbf{efficienza per riflessione}.

\section{Obiettivi e struttura della ricerca}
\label{sec:obiettivi}

La motivazione profonda di questa tesi risiede dunque nell'analisi di questo dualismo tra \textbf{imitazione strutturale} e \textbf{riflessione procedurale}. Si intende indagare se la strada per un'intelligenza artificiale realmente integrabile nella quotidianità passi per la riduzione dei modelli esistenti tramite distillazione o per l'elevazione delle loro capacità critiche tramite meccanismi di auto-correzione, o ancora — ipotesi più ambiziosa — se le due strategie possano rivelarsi complementari e sinergiche.

L'obiettivo principale del lavoro è quello di analizzare e mettere a confronto, attraverso un'indagine teorica e una valutazione delle performance riportate in letteratura, queste due tecniche. La ricerca si propone di:
\begin{itemize}
    \item Ricostruire i fondamenti teorici della Knowledge Distillation, analizzando come l'informazione venga compressa e trasferita senza perdite catastrofiche di accuratezza, dalle sue origini nei lavori di Hinton et al. alle varianti più recenti applicate agli LLM.
    \item Analizzare le architetture di feedback iterativo proprie del Reflection Tuning, valutando come l'introduzione di \textit{token di pensiero} e di cicli espliciti di auto-revisione modifichi la natura del ragionamento artificiale e riduca la frequenza delle allucinazioni.
    \item Identificare i contesti applicativi in cui una tecnica prevale sull'altra, valutando il trade-off tra velocità computazionale, sostenibilità infrastrutturale e solidità logica degli output.
    \item Esaminare il potenziale sinergico di questi due approcci, ipotizzando e analizzando scenari in cui modelli distillati possano essere ulteriormente potenziati tramite fasi di addestramento riflessivo, ottenendo sistemi che siano al contempo leggeri e logicamente affidabili.
\end{itemize}

In ultima analisi, il lavoro intende fornire una bussola metodologica per orientarsi tra le diverse strategie di ottimizzazione, contribuendo al dibattito su come sviluppare sistemi di Machine Learning che siano, al contempo, sostenibili, accessibili e logicamente affidabili per l'utente finale.


%*******************************************************************************
%*********************************** Parte rimossa *****************************
%*******************************************************************************

\begin{comment} \section{Struttura della tesi}

La tesi è articolata in cinque capitoli, ciascuno dedicato a un aspetto specifico dell’analisi e del confronto tra Knowledge Distillation e Reflection Tuning.

Il \hyperref[cap:1]{\textbf{Capitolo 1}} introduce il contesto di riferimento, le motivazioni che hanno guidato la scelta dell’argomento e gli obiettivi del lavoro.

Il \hyperref[cap:2]{\textbf{Capitolo 2}} è dedicato alle tecnologie utilizzate. In particolare, vengono richiamati i concetti fondamentali del Machine Learning e del Deep Learning, con un approfondimento sui Large Language Models e sull’architettura Transformer. Il capitolo prosegue con la descrizione delle tecniche di Knowledge Distillation e Reflection Tuning, analizzandone le principali definizioni, varianti, vantaggi e limiti, e si conclude con un confronto concettuale preliminare tra i due approcci.

Il \hyperref[cap:3]{\textbf{Capitolo 3}} descrive lo scenario applicativo, presentando le pipeline adottate per la Knowledge Distillation e per il Reflection Tuning, insieme al setup sperimentale e alla metodologia utilizzata per il confronto.

Nel \hyperref[cap:4]{\textbf{Capitolo 4}} vengono presentati e analizzati i risultati ottenuti. Il capitolo include una valutazione separata delle due tecniche, seguita da un’analisi comparativa e da una discussione critica dei risultati.

Infine, il \hyperref[cap:5]{\textbf{Capitolo 5}} riporta le conclusioni del lavoro, riassumendo i principali contributi emersi e delineando possibili sviluppi futuri.\end{comment}
